% Draft subsection for the Methods section of main.tex.
% Suggested placement: after \subsection{Synthetic Facial Image Generation},
% before \subsection{Landmark visualization}.
%
% Scope note: following the style of the other Methods subsections, this text
%   describes only how the assessment is carried out. The rationale for
%   comparing against a held-out baseline instead of a fixed verification
%   threshold, together with all reported values, belongs to Results 3.3 and is
%   deliberately absent here.
%
% TODO before merging into main.tex:
%   - add .bib entries and replace the citation keys arcface / adaface / lvface
%     / insightface / lpips / nnaa with the real ones
%   - CONFIRM that the held-out split is patient-level disjoint from the
%     training split; the sentence asserting this is the single most
%     load-bearing claim of the whole privacy analysis
%   - fill <UNVERIFIED> bootstrap replicate count (config default is 10, which
%     is a smoke-test value; state the number actually used)
%   - confirm the LPIPS backbone actually used (config default: AlexNet)
%   - confirm the AdaFace and LVFace checkpoints (ir_50 trained on MS1MV2, and
%     LVFace-L_Glint360K, respectively)
%   - the paper contains no display equations; adversarial accuracy is
%     therefore defined in prose. If a reviewer asks for a formal definition,
%     add a single display equation here rather than several
%   - cohort counts, image counts and detection-failure rates are reported in
%     Results 3.3, not here

\subsection{Privacy Assessment of the Pre-trained Generators}

Because the generators distributed with FaceKit are trained on photographs of
real individuals, the synthetic images they produce are assessed before
release for residual similarity to the training data. Two aspects are examined
separately: whether a synthetic image reproduces the identity of a training
individual, and whether it reproduces the appearance of a training photograph.

For the identity assessment, the prepared images of each cohort are divided
into a training partition, which is used to fit the generator, and a held-out
partition, which is disjoint from it at the patient level and never enters
generator training. Synthetic images are then sampled from the corresponding
generator. Faces are detected with InsightFace~\cite{insightface}, and the
five detected keypoints are used to align each face to a $112\times112$ crop
through a similarity transform. Identity embeddings are extracted with
ArcFace~\cite{arcface}, and the entire procedure is repeated with
AdaFace~\cite{adaface} and LVFace~\cite{lvface} so that the assessment does not
depend on a single recognition model; images in which no face is detected are
excluded. Distances between embeddings are measured as cosine distances, and
every synthetic and every held-out image is characterized by its distance to
the nearest training image of the same cohort. Rather than fixing a
verification threshold in advance, the threshold is read off the held-out
distribution: for a given percentile $p$, it is set to the $p$-th percentile of
the held-out nearest-neighbour distances, so that the operating point
corresponds to a fixed false-acceptance rate among real images of that cohort.
The proportion of synthetic images falling below this threshold is recorded at
$p = 1$, $2$, $5$, $10$ and $20$. The ordinal behaviour of each embedding model
is verified on the same images using triplets in which the positive pair
consists of two images of one patient acquired less than half a year apart and
the negative is an image of a different patient from the same cohort, checking
in each triplet whether the positive pair is the closer of the two.

The appearance assessment follows the same design with a perceptual rather
than an identity metric. LPIPS distances~\cite{lpips} are computed on images
resized to $256\times256$, and each synthetic and held-out image is compared
exhaustively against every training image of its cohort, retaining the
smallest distance; thresholds are again taken from the percentiles of the
held-out distribution. Both assessments are finally summarized by the nearest
neighbour adversarial accuracy~\cite{nnaa}. For two sets of images, the
adversarial accuracy is the proportion of images whose nearest neighbour in
the other set is farther away than their nearest neighbour within their own
set, averaged over the two directions, so that a value of one half indicates
that the two sets cannot be distinguished by proximity. It is evaluated
between the synthetic images and the training partition, and between the
synthetic images and the held-out partition, and the difference between the
two, termed the privacy loss, expresses whether synthetic images lie closer to
the training data than held-out real images of the same cohorts do. The
quantity is computed with embedding distances and with LPIPS distances, pooled
across cohorts, and confidence intervals are obtained by resampling images
with replacement over <UNVERIFIED> bootstrap replicates.
